% References use stable keys; numbered citations are generated by LaTeX.
\begin{thebibliography}{99}
\bibitem{ref006}
Edler, J., Fier, H. \& Grimpe, C. International scientist mobility and the locus of knowledge and technology transfer. \textit{Res. Policy} \textbf{40}, 791--805 (2011). \url{https://doi.org/10.1016/j.respol.2011.03.003}

\bibitem{ref009}
Chinchilla-Rodríguez, Z. et al. A global comparison of scientific mobility and collaboration according to national scientific capacities. \textit{Front. Res. Metr. Anal.} \textbf{3}, 17 (2018). \url{https://doi.org/10.3389/frma.2018.00017}

\bibitem{ref004}
Robinson-Garcia, N. et al. The many faces of mobility: Using bibliometric data to measure the movement of scientists. \textit{J. Informetr.} \textbf{13}, 50--63 (2019). \url{https://doi.org/10.1016/j.joi.2018.11.002}

\bibitem{ref022}
Moed, H. F. \& Halevi, G. A bibliometric approach to tracking international scientific migration. \textit{Scientometrics} \textbf{101}, 1987--2001 (2014). \url{https://doi.org/10.1007/s11192-014-1307-6}

\bibitem{ref012}
Sanliturk, E., Zagheni, E., Dańko, M. J., Theile, T. \& Akbaritabar, A. Global patterns of migration of scholars with economic development. \textit{Proc. Natl Acad. Sci. USA} \textbf{120}, e2217937120 (2023). \url{https://doi.org/10.1073/pnas.2217937120}

\bibitem{ref013}
Akbaritabar, A., Theile, T. \& Zagheni, E. Bilateral flows and rates of international migration of scholars for 210 countries for the period 1998--2020. \textit{Sci. Data} \textbf{11}, 816 (2024). \url{https://doi.org/10.1038/s41597-024-03655-9}

\bibitem{ref038}
Gomez, C. J., Herman, A. C. \& Parigi, P. Moving more, but closer: Mapping the growing regionalization of global scientific mobility using ORCID. \textit{J. Informetr.} \textbf{14}, 101044 (2020). \url{https://doi.org/10.1016/j.joi.2020.101044}

\bibitem{ref031}
Verginer, L. \& Riccaboni, M. Talent goes to global cities: The world network of scientists’ mobility. \textit{Res. Policy} \textbf{50}, 104127 (2021). \url{https://doi.org/10.1016/j.respol.2020.104127}

\bibitem{ref032}
Na, Y. \& Liu, X. Unveiling the core-periphery dynamics and driving factors of intercity scientific mobility in China. \textit{Habitat Int.} \textbf{163}, 103481 (2025). \url{https://doi.org/10.1016/j.habitatint.2025.103481}

\bibitem{ref029}
Akbaritabar, A., Dańko, M. J., Zhao, X. \& Zagheni, E. Global subnational estimates of migration of scientists reveal large disparities in internal and international flows. \textit{Proc. Natl Acad. Sci. USA} \textbf{122}, e2424521122 (2025). \url{https://doi.org/10.1073/pnas.2424521122}

\bibitem{ref033}
Momeni, F., Karimi, F., Mayr, P., Peters, I. \& Dietze, S. The many facets of academic mobility and its impact on scholars’ career. \textit{J. Informetr.} \textbf{16}, 101280 (2022). \url{https://doi.org/10.1016/j.joi.2022.101280}

\bibitem{ref034}
Subbotin, A. \& Aref, S. Brain drain and brain gain in Russia: Analyzing international migration of researchers by discipline using Scopus bibliometric data 1996--2020. \textit{Scientometrics} \textbf{126}, 7875--7900 (2021). \url{https://doi.org/10.1007/s11192-021-04091-x}

\bibitem{ref035}
Kim, H. Global mobility of the recent STEM postdoctoral workforce registered in ORCID. \textit{Quant. Sci. Stud.} \textbf{6}, 119--130 (2025). \url{https://doi.org/10.1162/qss_a_00341}

\bibitem{ref021}
Priem, J., Piwowar, H. \& Orr, R. OpenAlex: A fully-open index of scholarly works, authors, venues, institutions, and concepts. Preprint at arXiv:2205.01833 (2022). \url{https://doi.org/10.48550/arXiv.2205.01833}

\bibitem{ref037}
Czaika, M. \& de Haas, H. The globalization of migration: Has the world become more migratory? \textit{Int. Migr. Rev.} \textbf{48}, 283--323 (2014). \url{https://doi.org/10.1111/imre.12095}

\bibitem{ref039}
Zhao, X., Akbaritabar, A., Kashyap, R. \& Zagheni, E. A gender perspective on the global migration of scholars. \textit{Proc. Natl Acad. Sci. USA} \textbf{120}, e2214664120 (2023). \url{https://doi.org/10.1073/pnas.2214664120}

\bibitem{ref026}
Verginer, L. \& Riccaboni, M. Cities and countries in the global scientist mobility network. \textit{Appl. Netw. Sci.} \textbf{5}, 38 (2020). \url{https://doi.org/10.1007/s41109-020-00276-0}

\bibitem{ref043}
Nunes, A., Trappenberg, T. \& Alda, M. The definition and measurement of heterogeneity. \textit{Transl. Psychiatry} \textbf{10}, 299 (2020). \url{https://doi.org/10.1038/s41398-020-00986-0}

\bibitem{gotelli2001}
Gotelli, N. J. \& Colwell, R. K. Quantifying biodiversity: procedures and pitfalls in the measurement and comparison of species richness. \textit{Ecol. Lett.} \textbf{4}, 379--391 (2001). \url{https://doi.org/10.1046/j.1461-0248.2001.00230.x}

\bibitem{lenormand2016}
Lenormand, M., Bassolas, A. \& Ramasco, J. J. Systematic comparison of trip distribution laws and models. \textit{J. Transp. Geogr.} \textbf{51}, 158--169 (2016). \url{https://doi.org/10.1016/j.jtrangeo.2015.12.008}

\bibitem{ref027}
Vaccario, G., Verginer, L. \& Schweitzer, F. The mobility network of scientists: Analyzing temporal correlations in scientific careers. \textit{Appl. Netw. Sci.} \textbf{5}, 36 (2020). \url{https://doi.org/10.1007/s41109-020-00279-x}

\bibitem{zhou2007}
Zhou, T., Ren, J., Medo, M. \& Zhang, Y.-C. Bipartite network projection and personal recommendation. \textit{Phys. Rev. E} \textbf{76}, 046115 (2007). \url{https://doi.org/10.1103/PhysRevE.76.046115}

\bibitem{gould1989}
Gould, R. V. \& Fernandez, R. M. Structures of mediation: a formal approach to brokerage in transaction networks. \textit{Sociol. Methodol.} \textbf{19}, 89--126 (1989). \url{https://doi.org/10.2307/270949}

\bibitem{ref001}
Schott, T. World science: Globalization of institutions and participation. \textit{Sci. Technol. Hum. Values} \textbf{18}, 196--208 (1993). \url{https://doi.org/10.1177/016224399301800205}

\bibitem{ref002}
Marginson, S. What drives global science? The four competing narratives. \textit{Stud. High. Educ.} \textbf{47}, 1566--1584 (2022). \url{https://doi.org/10.1080/03075079.2021.1942822}

\bibitem{ref008}
Czaika, M. \& Orazbayev, S. The globalisation of scientific mobility, 1970--2014. \textit{Appl. Geogr.} \textbf{96}, 1--10 (2018). \url{https://doi.org/10.1016/j.apgeog.2018.04.017}

\bibitem{ref020}
Gomez, C. J. The growing concentration of national influence in global science. \textit{Nat. Hum. Behav.} \textbf{10}, 1645--1661 (2026). \url{https://doi.org/10.1038/s41562-026-02489-2}

\bibitem{ref041}
Docquier, F. \& Rapoport, H. Globalization, brain drain, and development. \textit{J. Econ. Lit.} \textbf{50}, 681--730 (2012). \url{https://doi.org/10.1257/jel.50.3.681}

\bibitem{ref042}
Saxenian, A. From brain drain to brain circulation: Transnational communities and regional upgrading in India and China. \textit{Stud. Comp. Int. Dev.} \textbf{40}, 35--61 (2005). \url{https://doi.org/10.1007/BF02686293}

\bibitem{narayan2022}
Narayan, A., Chami, I., Orr, L. \& R{\'e}, C. Can foundation models wrangle your data? \textit{Proc. VLDB Endow.} \textbf{16}, 738--746 (2022). \url{https://doi.org/10.14778/3574245.3574258}

\bibitem{peeters2025}
Peeters, R., Steiner, A. \& Bizer, C. Entity matching using large language models. In \textit{Proceedings of the 28th International Conference on Extending Database Technology} 529--541 (OpenProceedings.org, 2025). \url{https://doi.org/10.48786/edbt.2025.42}

\bibitem{rhoades1993}
Rhoades, S. A. The Herfindahl-Hirschman index. \textit{Fed. Reserve Bull.} \textbf{79}, 188--189 (1993). \url{https://fedinprint.org/item/fedgrb/15160}

\bibitem{hill1973}
Hill, M. O. Diversity and evenness: a unifying notation and its consequences. \textit{Ecology} \textbf{54}, 427--432 (1973). \url{https://doi.org/10.2307/1934352}

\bibitem{gini1921}
Gini, C. Measurement of inequality of incomes. \textit{Econ. J.} \textbf{31}, 124--126 (1921). \url{https://doi.org/10.2307/2223319}

\bibitem{baker1994}
Baker, S. G. The multinomial-Poisson transformation. \textit{J. R. Stat. Soc. D} \textbf{43}, 495--504 (1994). \url{https://doi.org/10.2307/2348134}

\bibitem{byrd1995}
Byrd, R. H., Lu, P., Nocedal, J. \& Zhu, C. A limited memory algorithm for bound constrained optimization. \textit{SIAM J. Sci. Comput.} \textbf{16}, 1190--1208 (1995). \url{https://doi.org/10.1137/0916069}

\bibitem{brin1998}
Brin, S. \& Page, L. The anatomy of a large-scale hypertextual Web search engine. \textit{Comput. Netw. ISDN Syst.} \textbf{30}, 107--117 (1998). \url{https://doi.org/10.1016/S0169-7552(98)00110-X}

\bibitem{salton1988}
Salton, G. \& Buckley, C. Term-weighting approaches in automatic text retrieval. \textit{Inf. Process. Manag.} \textbf{24}, 513--523 (1988). \url{https://doi.org/10.1016/0306-4573(88)90021-0}

\bibitem{nelder1972}
Nelder, J. A. \& Wedderburn, R. W. M. Generalized linear models. \textit{J. R. Stat. Soc. A} \textbf{135}, 370--384 (1972). \url{https://doi.org/10.2307/2344614}

\bibitem{kvalseth1985}
Kvålseth, T. O. Cautionary note about $R^2$. \textit{Am. Stat.} \textbf{39}, 279--285 (1985). \url{https://doi.org/10.1080/00031305.1985.10479448}

\bibitem{hodson2022}
Hodson, T. O. Root-mean-square error (RMSE) or mean absolute error (MAE): when to use them or not. \textit{Geosci. Model Dev.} \textbf{15}, 5481--5487 (2022). \url{https://doi.org/10.5194/gmd-15-5481-2022}

\bibitem{spearman1904}
Spearman, C. The proof and measurement of association between two things. \textit{Am. J. Psychol.} \textbf{15}, 72--101 (1904). \url{https://doi.org/10.2307/1412159}

\bibitem{long2000}
Long, J. S. \& Ervin, L. H. Using heteroscedasticity consistent standard errors in the linear regression model. \textit{Am. Stat.} \textbf{54}, 217--224 (2000). \url{https://doi.org/10.1080/00031305.2000.10474549}

\bibitem{efron1979}
Efron, B. Bootstrap methods: another look at the jackknife. \textit{Ann. Stat.} \textbf{7}, 1--26 (1979). \url{https://doi.org/10.1214/aos/1176344552}

\bibitem{holm1979}
Holm, S. A simple sequentially rejective multiple test procedure. \textit{Scand. J. Stat.} \textbf{6}, 65--70 (1979). \url{https://www.jstor.org/stable/4615733}

\bibitem{cameron2015}
Cameron, A. C. \& Miller, D. L. A practitioner's guide to cluster-robust inference. \textit{J. Hum. Resour.} \textbf{50}, 317--372 (2015). \url{https://doi.org/10.3368/jhr.50.2.317}

\bibitem{ref003}
Sugimoto, C. R. et al. Scientists have most impact when they're free to move. \textit{Nature} \textbf{550}, 29--31 (2017). \url{https://doi.org/10.1038/550029a}

\bibitem{ref005}
Deville, P. et al. Career on the move: Geography, stratification and scientific impact. \textit{Sci. Rep.} \textbf{4}, 4770 (2014). \url{https://doi.org/10.1038/srep04770}

\bibitem{ref007}
Netz, N., Hampel, S. \& Aman, V. What effects does international mobility have on scientists’ careers? A systematic review. \textit{Res. Eval.} \textbf{29}, 327--351 (2020). \url{https://doi.org/10.1093/reseval/rvaa007}

\bibitem{ref010}
Appelt, S., van Beuzekom, B., Galindo-Rueda, F. \& de Pinho, R. Which factors influence the international mobility of research scientists? OECD Science, Technology and Industry Working Papers 2015/02 (OECD Publishing, 2015). \url{https://doi.org/10.1787/5js1tmrr2233-en}

\bibitem{ref011}
Franzoni, C., Scellato, G. \& Stephan, P. Foreign-born scientists: Mobility patterns for 16 countries. \textit{Nat. Biotechnol.} \textbf{30}, 1250--1253 (2012). \url{https://doi.org/10.1038/nbt.2449}

\bibitem{ref014}
Leydesdorff, L. \& Wagner, C. S. International collaboration in science and the formation of a core group. \textit{J. Informetr.} \textbf{2}, 317--325 (2008). \url{https://doi.org/10.1016/j.joi.2008.07.003}

\bibitem{ref015}
Grossetti, M. et al. Cities and the geographical deconcentration of scientific activity: A multilevel analysis of publications (1987--2007). \textit{Urban Stud.} \textbf{51}, 2219--2234 (2014). \url{https://doi.org/10.1177/0042098013506047}

\bibitem{ref016}
Maisonobe, M., Eckert, D., Grossetti, M., Jégou, L. \& Milard, B. The world network of scientific collaborations between cities: Domestic or international dynamics? \textit{J. Informetr.} \textbf{10}, 1025--1036 (2016). \url{https://doi.org/10.1016/j.joi.2016.06.002}

\bibitem{ref017}
Kitajima, K. \& Okamura, K. The altering landscape of US--China science collaboration: From convergence to divergence. \textit{Humanit. Soc. Sci. Commun.} \textbf{12}, 302 (2025). \url{https://doi.org/10.1057/s41599-025-04550-3}

\bibitem{ref018}
Nielsen, M. W. \& Andersen, J. P. Global citation inequality is on the rise. \textit{Proc. Natl Acad. Sci. USA} \textbf{118}, e2012208118 (2021). \url{https://doi.org/10.1073/pnas.2012208118}

\bibitem{ref019}
Gomez, C. J., Herman, A. C. \& Parigi, P. Leading countries in global science increasingly receive more citations than other countries doing similar research. \textit{Nat. Hum. Behav.} \textbf{6}, 919--929 (2022). \url{https://doi.org/10.1038/s41562-022-01351-5}

\bibitem{ref023}
Hoekman, J., Frenken, K. \& Tijssen, R. J. W. Research collaboration at a distance: Changing spatial patterns of scientific collaboration within Europe. \textit{Res. Policy} \textbf{39}, 662--673 (2010). \url{https://doi.org/10.1016/j.respol.2010.01.012}

\bibitem{ref024}
Pan, R., Kaski, K. \& Fortunato, S. World citation and collaboration networks: Uncovering the role of geography in science. \textit{Sci. Rep.} \textbf{2}, 902 (2012). \url{https://doi.org/10.1038/srep00902}

\bibitem{ref025}
Hoekman, J., Frenken, K. \& van Oort, F. The geography of collaborative knowledge production in Europe. \textit{Ann. Reg. Sci.} \textbf{43}, 721--738 (2009). \url{https://doi.org/10.1007/s00168-008-0252-9}

\bibitem{ref028}
Vaccario, G., Verginer, L. \& Schweitzer, F. Reproducing scientists’ mobility: A data-driven model. \textit{Sci. Rep.} \textbf{11}, 10733 (2021). \url{https://doi.org/10.1038/s41598-021-90281-9}

\bibitem{ref030}
Slocombe, J., Rowe, F. \& Akbaritabar, A. Global sequencing of researcher migration reveals factors associated with disparities in subnational patterns of brain gain and drain. MPIDR Working Paper WP-2026-018 (2026). \url{https://doi.org/10.4054/MPIDR-WP-2026-018}

\bibitem{ref036}
Gureyev, V. N., Mazov, N. A., Kosyakov, D. V. \& Guskov, A. E. Review and analysis of publications on scientific mobility: Assessment of influence, motivation, and trends. \textit{Scientometrics} \textbf{124}, 1599--1630 (2020). \url{https://doi.org/10.1007/s11192-020-03515-4}

\bibitem{ref040}
Cao, C., Baas, J., Wagner, C. S. \& Jonkers, K. Returning scientists and the emergence of China’s science system. \textit{Sci. Public Policy} \textbf{47}, 172--183 (2020). \url{https://doi.org/10.1093/scipol/scz056}

\bibitem{ref044}
Barrat, A., Barthélemy, M., Pastor-Satorras, R. \& Vespignani, A. The architecture of complex weighted networks. \textit{Proc. Natl Acad. Sci. USA} \textbf{101}, 3747--3752 (2004). \url{https://doi.org/10.1073/pnas.0400087101}

\bibitem{ref045}
Huang, J., Cao, C. \& Liu, H. Indigenization and inclusion in Chinese academia. \textit{Nat. Hum. Behav.} \textbf{10}, 492--504 (2026). \url{https://doi.org/10.1038/s41562-025-02366-4}

\bibitem{ref046}
Wang, H., Yang, G. \& Yue, Z. Breaking through ingrained beliefs: Revisiting the impact of the digital economy on carbon emissions. \textit{Humanit. Soc. Sci. Commun.} \textbf{10}, 609 (2023). \url{https://doi.org/10.1057/s41599-023-02126-7}

\bibitem{yenilmez2026}
Yenilmez, T. Fifty years of the world academic migration network. \textit{Appl. Netw. Sci.} \textbf{11}, 65 (2026). \url{https://doi.org/10.1007/s41109-026-00840-0}

\end{thebibliography}
